(1) 对围绕无限直导线1的半径 \(d\) 的圆形安培回路，\(\oint\mathbf B\cdot d\boldsymbol\ell=2\pi d B=\mu_0I_1\)。因此导线2位置的磁场大小为 \(B_1=\mu_0|I_1|/(2\pi d)\)，方向沿环绕导线1的切向。
(2) 导线2的电流元受力 \(d\mathbf F=I_2d\boldsymbol\ell\times\mathbf B_1\)，两向量互相垂直，故单位长度力大小 \(f=\mu_0|I_1I_2|/(2\pi d)\)。力沿两线的连心方向，与导线长度无关是无限均匀直线近似的结果。
(3) 当两电流同向，右手定则与叉乘给相互吸引；电流反向时则相互排斥。交换导线1、2，力大小因 \(I_1I_2\) 对称不变，连线的指向恰好反号，因此两导线所受力大小相等、方向相反，符合稳恒机械动量平衡。
(4) \(\mu_0\) 量纲为牛每安培平方，故 \(\mu_0I^2/d\) 单位为牛每米。有限导线在端部存在场线弯曲，交流电流涉及传播延迟、感应与辐射，未必可用同一时刻的无限直流表达式求全系统的即时作用力。
